% CERT-SCAN preprint. Numbers come from stats.tex, generated by code/07_preprint.py.
\documentclass[10pt,twocolumn,letterpaper]{article}
\newcommand{\nJuris}{52}
\newcommand{\nRows}{104}
\newcommand{\nReqJuris}{2}
\newcommand{\nReqJurisW}{two}
\newcommand{\nCredJuris}{5}
\newcommand{\nCredJurisW}{five}
\newcommand{\nGenJuris}{1}
\newcommand{\nGenJurisW}{one}
\newcommand{\nAnyJuris}{8}
\newcommand{\nAnyJurisW}{eight}
\newcommand{\nNoneJuris}{44}
\newcommand{\dwReq}{2}
\newcommand{\dwCred}{4}
\newcommand{\dwGen}{1}
\newcommand{\dwNone}{45}
\newcommand{\wwReq}{2}
\newcommand{\wwCred}{3}
\newcommand{\wwGen}{1}
\newcommand{\wwNone}{46}
\newcommand{\nWPIdw}{33}
\newcommand{\nWPIww}{34}
\newcommand{\nOtherExamDW}{18}
\newcommand{\nOtherExamWW}{18}
\newcommand{\nHigh}{8}
\newcommand{\nMedium}{94}
\newcommand{\nLow}{2}
\newcommand{\nVerified}{12}
\newcommand{\nStateSpecific}{102}
\newcommand{\nQA}{0}
\newcommand{\nPositive}{13}
\newcommand{\nInstStatute}{2}
\newcommand{\nInstReg}{4}
\newcommand{\nInstGuidance}{6}
\newcommand{\nInstBoard}{1}
\newcommand{\epaPages}{131}
\newcommand{\epaCyber}{0}
\newcommand{\epaSecurity}{1}
\newcommand{\epaScada}{1}
\newcommand{\nReqJurisCap}{Two}
\newcommand{\nCredJurisCap}{Five}
\newcommand{\reqList}{Maryland and New York}
\newcommand{\credList}{Maine, Massachusetts, Michigan, Montana, and Texas}
\newcommand{\nWPIdwNoReq}{31}
\newcommand{\nNonCodified}{7}
\newcommand{\accessDate}{23 October 2026}
\newcommand{\authorName}{Deborah Wale Ibinayo}
\newcommand{\authorOrcid}{0009-0002-3571-5981}
\newcommand{\authorEmail}{deborahwaleibinayo@gmail.com}
\newcommand{\authorAffil}{Independent Researcher, Lawrenceville, Georgia, USA}
\newcommand{\coName}{Dayo Oladipo}
\newcommand{\coOrcid}{0000-0001-7724-2903}
\newcommand{\coEmail}{oladipodayogeorge@gmail.com}
\newcommand{\coRoles}{[roles to be confirmed]}
\newcommand{\repoURL}{[GitHub URL assigned at release]}
\newcommand{\doiText}{[DOI assigned at release]}
\newif\iffinal
\finalfalse

\newcommand{\tabRequired}{Maryland & DW & New and renewal training & 0.5 & 3 & 2025-10-01 & Statute & Md. Code Envir. 9-2702(2) (2025 Md. Laws ch. 495; SB 871) \\
Maryland & WW & New and renewal training & 0.5 & 3 & 2025-10-01 & Statute & Md. Code Envir. 9-2702(2) (2025 Md. Laws ch. 495; SB 871) \\
New York & DW & New and renewal training & 1 & 3 & 2026-03-11 & Regulation & 10 NYCRR Part 5 Appendix 5-E section 5-E.7 \\
New York & WW & Renewal hours (within CE total) & 2-4 by grade & 5 & 2026-03-26 & Regulation & 6 NYCRR 650.8 and 650.12 \\
}
\newcommand{\tabCredited}{Maine & WW & Cyber eligible for CE credit & Agency guidance & 2022 & Maine DEP, TCH Requirements for Certification Renewal v1.3 (2022), Table 1 \\
Massachusetts & DW & Cyber eligible for CE credit & Board course approvals & 2026-03-13 & Board TCH approvals DWT-2026-05 and DWT-2026-09 to -13; MassDEP self-paced cybersecurity course (1 TCH) \\
Michigan & DW & Cyber eligible for CE credit & Agency guidance & 2024-06 & EGLE approved course code 4274 (0.2 technical DW CEC) \\
Michigan & WW & Cyber eligible for CE credit & Agency guidance & 2024-06 & EGLE approved course code 3315 (0.3 managerial WW CEC) \\
Montana & DW & Cyber eligible for CE credit & Agency guidance & 2023-06-08 & Montana DEQ Continuing Education Credit Guidelines (2023-06-08), sections 3.1 and 5.2 \\
Montana & WW & Cyber eligible for CE credit & Agency guidance & 2023-06-08 & Montana DEQ Continuing Education Credit Guidelines (2023-06-08), sections 3.1 and 5.2 \\
Texas & DW & Cyber optional within mandatory resiliency hours & Agency guidance & 2024-04-01 & TCEQ RG-637 Resiliency Requirements for Water Operators and Training Providers (30 TAC ch. 30 section not identified) \\
}
\newcommand{\tabExam}{Water treatment & Listed & Listed & Listed & Listed \\
Water distribution & -- & -- & -- & -- \\
Wastewater treatment & -- & -- & -- & -- \\
Wastewater collection & -- & -- & -- & -- \\
}
\newcommand{\tabAppendix}{Alabama & -- & -- & WPI+ & WPI & M & Ala. Code 22-25-1 to -15; ADEM Admin. Code r. 335-10-1 \\
Alaska & -- & -- & WPI+ & WPI & M & AS 46.30; 18 AAC 74 \\
Arizona & -- & -- & WPI+ & WPI & M & ARS 49-351; A.A.C. R18-5 \\
Arkansas & -- & -- & WPI+ & WPI & M & Ark. Code 17-51; ADH Rules; WW: Ark. Code 8-5; APC\&EC Rule 3 \\
California & -- & -- & Other & Other & M & Cal. H\&S Code 106875-106910; 22 CCR 63750 et seq.; WW: Cal. Water Code 13625 et seq.; 23 CCR 3670 et seq. \\
Colorado & -- & -- & WPI+ & WPI & M & CRS 25-9; 5 CCR 1003-2 (Reg. 100) \\
Connecticut & -- & -- & WPI+ & WPI & M & Conn. Gen. Stat. 25-32n; RCSA 25-32-7a to -14; WW: Conn. Gen. Stat. 22a-416 et seq.; RCSA 22a-416-1 et seq. \\
Delaware & -- & -- & Other & WPI & M & 16 Del. C. 122(3)(c); 4 DE Admin. Code 4462; WW: 7 Del. C. ch. 60; 7 DE Admin. Code 7204 \\
District of Columbia & -- & -- & Other & Other & L & No independent District program identified; confirmation requested \\
Florida & -- & -- & Other & Other & H & F.S. 403.865-.876; F.A.C. 62-602 \\
Georgia & -- & -- & WPI+ & WPI & M & OCGA 43-51; Ga. Comp. R. \& Regs. 750 \\
Hawaii & -- & -- & WPI+ & WPI & M & HRS 340F; HAR 11-25; WW: HRS 340B; HAR 11-61 \\
Idaho & -- & -- & WPI+ & WPI & M & Idaho Code 54-24; IDAPA 24.05.01 \\
Illinois & -- & -- & Other & Other & M & 415 ILCS 45; 35 Ill. Adm. Code 681; WW: 415 ILCS 5/13; 35 Ill. Adm. Code 380 \\
Indiana & -- & -- & Other & Other & H/M & IC 13-18-11; 327 IAC 8-12; WW: IC 13-18-11; 327 IAC 5-22 \\
Iowa & -- & -- & WPI+ & WPI & M & Iowa Code 455B.211-.224; 567 IAC 81 \\
Kansas & -- & -- & WPI+ & WPI & M & KSA 65-4501 et seq.; KAR 28-16-28 to -35 \\
Kentucky & -- & -- & Other & Other & M & KRS 223 and 224; 401 KAR 11 \\
Louisiana & -- & -- & Other & Other & M & La. R.S. 40:1141-1151; LAC 48:V ch. 73 \\
Maine & -- & C & WPI+ & WPI & M & 32 MRSA 4171 et seq.; 10-144 CMR 90-429; WW: 32 MRSA 4171; 06-096 CMR 531 \\
Maryland & R & R & WPI+ & WPI & H & Md. Code Envir. Title 12; COMAR 26.06.01 \\
Massachusetts & C & -- & WPI+ & WPI & M & MGL c.13 s.66B; 236 CMR; WW: MGL c.13 s.80-81; 257 CMR 2 \\
Michigan & C & C & Other & Other & M & MCL 325.1009; Mich. Admin. Code R 325.11901 et seq.; WW: Part 31 NREPA; Mich. Admin. Code R 299.2901 et seq. \\
Minnesota & -- & -- & WPI+ & Other & M & Minn. Stat. 115.71-.77; Minn. R. 9400 \\
Mississippi & -- & -- & WPI+ & WPI & M & Miss. Code 21-27-201 et seq.; 15 Miss. Admin. Code Pt. 20; WW: Miss. Code 21-27-201 et seq. \\
Missouri & -- & -- & Other & WPI & M & RSMo 640.100; 10 CSR 60-14 and 20 CSR 2050 \\
Montana & C & C & WPI+ & WPI & M & MCA 37-42; ARM 17.40.2 \\
Nebraska & -- & -- & Other & WPI & M & Neb. Rev. Stat. 71-5301 et seq.; 179 NAC 10; WW: Neb. Rev. Stat. 81-15; Title 197 \\
Nevada & -- & -- & WPI+ & Other & M & NRS 445A.875-.880; NAC 445A.617-.652; WW: NRS 445A; NAC 445A.285-.2912 \\
New Hampshire & -- & -- & WPI+ & Other & M & RSA 332-E; Env-Dw 502; WW: RSA 485-A:5; Env-Wq 900 \\
New Jersey & -- & -- & WPI+ & WPI & M & NJSA 58:11-64 to -73; N.J.A.C. 7:10A \\
New Mexico & -- & -- & WPI & WPI & M & NMSA 61-33; 20.7.4 NMAC \\
New York & R & R & WPI+ & WPI & H/M & 10 NYCRR Subpart 5-4; WW: 6 NYCRR Part 650 \\
North Carolina & -- & -- & Other & Other & M & GS 90A-20 to -32; 15A NCAC 18D; WW: GS 143B-300; 15A NCAC 08G \\
North Dakota & -- & -- & Other & Other & M & NDCC 23.1-; NDAC 33.1-19-01 \\
Ohio & -- & -- & WPI+ & WPI & H & ORC 6109; OAC 3745-7; WW: ORC 6111; OAC 3745-7 \\
Oklahoma & -- & -- & WPI+ & WPI & M & 27A O.S. 2-1-101; OAC 252:710 \\
Oregon & -- & -- & WPI+ & WPI & M & ORS 448.131; OAR 333-061; WW: ORS 448.405-.470; OAR 340-049 \\
Pennsylvania & G & G & Other & Other & M & 63 P.S. 1001 et seq.; 25 Pa. Code ch. 302 \\
Puerto Rico & -- & -- & Other & Other & M & Ley 53-1978; Reglamento 8750 \\
Rhode Island & -- & -- & WPI+ & Other & M & RIGL 23-65; 216-RICR-50-05-5; WW: RIGL 42-17.1; 250-RICR-150-10-2 \\
South Carolina & -- & -- & WPI+ & WPI & M & SC Code 40-23; SC Regs ch. 51 \\
South Dakota & -- & -- & WPI+ & WPI & M & SDCL 34A-3; ARSD 74:21 \\
Tennessee & -- & -- & WPI+ & WPI & M & TCA 68-221-901 to -915; Tenn. Comp. R. \& Regs. 0400-49-01 \\
Texas & C & -- & Other & Other & M & Tex. Water Code ch. 37; 30 TAC ch. 30 \\
Utah & -- & -- & Other & WPI & M & Utah Code 19-4-104; Utah Admin. Code R309-300; WW: Utah Code 19-5; Utah Admin. Code R317-10 \\
Vermont & -- & -- & WPI+ & WPI & M & 10 V.S.A. ch. 56; Water Supply Rule ch. 21; WW: 10 VSA ch. 47; Wastewater Treatment Facility Operator Certification Rule \\
Virginia & -- & -- & WPI+ & WPI & M & Va. Code 54.1-2300 et seq.; 18 VAC 160-30 \\
Washington & -- & -- & WPI+ & WPI & M & RCW 70.119; WAC 246-292; WW: RCW 70A.212; WAC 173-230 \\
West Virginia & -- & -- & WPI+ & WPI & M & W. Va. Code 16-1-4 and 16-1-9; 64 CSR 4; WW: W. Va. Code 16-1; 64 CSR 4 and 64 CSR 5 \\
Wisconsin & -- & -- & Other & Other & M & Wis. Stat. 281.17(3); Wis. Admin. Code NR 114 \\
Wyoming & -- & -- & WPI+ & WPI & M & WS 35-11-302; WQRR ch. 5 \\
}

\usepackage[T1]{fontenc}
\usepackage{mathpazo}
\usepackage[margin=0.72in,columnsep=0.28in,top=0.8in,bottom=0.8in]{geometry}
\usepackage{microtype}
\usepackage{graphicx}
\usepackage{booktabs,array,tabularx,longtable,multirow}
\usepackage{ragged2e}
\usepackage[hyphens]{xurl}
\usepackage[dvipsnames]{xcolor}
\usepackage[colorlinks=true,linkcolor=red!70!black,citecolor=blue!70!black,urlcolor=blue!60!black]{hyperref}
\usepackage{orcidlink}
\usepackage{titlesec}
\usepackage[font=small,labelsep=colon,justification=centering]{caption}
\usepackage{enumitem}
\usepackage{fancyhdr}
\usepackage{amssymb}

% ---- APS-like headings: I. INTRODUCTION / A. Subsection
\renewcommand{\thesection}{\Roman{section}}
\renewcommand{\thesubsection}{\Alph{subsection}}
\titleformat{\section}{\centering\normalfont\small\bfseries}{\thesection.}{0.6em}{\MakeUppercase}
\titlespacing*{\section}{0pt}{1.6ex plus .4ex}{1.0ex}
\titleformat{\subsection}{\centering\normalfont\small\bfseries}{\thesubsection.}{0.6em}{}
\titlespacing*{\subsection}{0pt}{1.3ex plus .3ex}{0.8ex}
\renewcommand{\figurename}{FIG.}
\renewcommand{\tablename}{TABLE}
\renewcommand{\thetable}{\Roman{table}}
\captionsetup[table]{position=top,skip=4pt}
\setlength{\parskip}{0pt}
\setlength{\parindent}{1.2em}
\setlist{nosep,leftmargin=1.3em}
\pagestyle{fancy}\fancyhf{}\fancyhead[R]{\small\thepage}\renewcommand{\headrulewidth}{0pt}
\fancypagestyle{plain}{\fancyhf{}\renewcommand{\headrulewidth}{0pt}}
\newcommand{\figw}{\columnwidth}

\begin{document}
\twocolumn[{%
\begin{center}
{\large\bfseries CERT-SCAN: Cybersecurity in U.S.\ Drinking Water and Wastewater\\[2pt] Operator Certification Across \nJuris\ Jurisdictions}\\[10pt]
{\normalsize \authorName\,\orcidlink{\authorOrcid}\,\textsuperscript{1,\,*} and \coName\,\orcidlink{\coOrcid}\,\textsuperscript{2,\,\dag}}\\[4pt]
{\small\itshape \textsuperscript{1}Independent Researcher, Lawrenceville, Georgia, USA}\\{\small\itshape \textsuperscript{2}Independent Researcher, Troy, New York, USA}\\[4pt]
\end{center}
\begin{minipage}{0.86\textwidth}\small
\setlength{\parindent}{1.2em}
\hspace{1.2em}Drinking water and wastewater systems in the United States are run by state-certified operators, and the control systems those operators use have become a recurring target of cyberattacks. Whether states require operators to know anything about cybersecurity to obtain or keep their certification has never been tracked nationally. The last national summary of state operator certification programs, published by the U.S.\ Environmental Protection Agency in March 2016, does not use the word ``cyber'' (\epaCyber\ occurrences in \epaPages\ pages). We present CERT-SCAN, an open, structured, and annually maintained record of how all 50 states, the District of Columbia, and Puerto Rico treat cybersecurity in operator certification, coded separately for drinking water and wastewater (\nRows\ records), with the statute, regulation, or agency document behind every record. As of \accessDate, \nReqJurisW\ jurisdictions require cybersecurity training by statute or regulation for both programs: Maryland, by statute effective 1 October 2025, and New York, by regulations adopted in March 2026. \nCredJurisCap\ jurisdictions (\credList) credit cybersecurity training toward renewal without requiring it, and Pennsylvania has required a one-time general security course since 2010 whose cybersecurity content is unconfirmed. In the remaining jurisdictions, no provision was found. CERT-SCAN also documents a gap between testing and training: the 2025 exam content outlines of Water Professionals International, used by \nWPIdw\ drinking water programs, list cybersecurity as a task for water treatment operators at every class level, while no distribution, wastewater treatment, or collection outline does. Most operators who may now be tested on cybersecurity work in states that do not require them to be trained in it. We describe the coding scheme, source protocol, validation, and limitations, and release the data, code, and figures openly.\\[4pt]
\textbf{Keywords:} water sector cybersecurity, operator certification, drinking water, wastewater, state regulation, workforce, open data
\end{minipage}
\vspace{12pt}
}]
{\renewcommand{\thefootnote}{\fnsymbol{footnote}}\footnotetext[1]{Corresponding author: \href{mailto:\authorEmail}{\authorEmail}}\footnotetext[2]{\href{mailto:\coEmail}{\coEmail}}}

\section{Introduction}

Every community water system and non-transient non-community water system in the United States must be run by an operator certified by the state. The requirement flows from Section 1419 of the Safe Drinking Water Act \cite{sdwa} and from the operator certification guidelines EPA published on 5 February 1999 \cite{epa1999}: a state that does not implement a program meeting the guidelines has 20 percent of its Drinking Water State Revolving Fund allotment withheld \cite{epaopcert}. Wastewater operators are certified under separate state authorities. Because certification reaches every operator at every regulated system, including the many small systems with no information technology staff, it is one of the few policy levers that can shape what the entire operating workforce knows.

The systems those operators control have changed. Pumps, valves, chemical feeds, and storage levels are now commonly managed through programmable logic controllers and supervisory control and data acquisition (SCADA) systems, some reachable from the internet. In November 2023, attackers affiliated with Iran's Islamic Revolutionary Guard Corps compromised a Unitronics controller at a booster station of the Municipal Water Authority of Aliquippa, Pennsylvania; the device was reportedly exposed to the internet and protected by a factory-default password \cite{securityweek}. In May 2024, EPA reported that more than 70 percent of the drinking water systems it had inspected since September 2023 did not meet basic Section 1433 requirements, and named practices such as unchanged default passwords and shared logins; its recommended actions included cybersecurity awareness training \cite{epa2024}. In the summer of 2026, coordinated attacks targeted water and wastewater utilities in at least twelve states, changing passwords on internet-facing controllers and locking out the employees who monitor them; some utilities switched to manual operation and issued precautionary boil-water notices \cite{axios2026}.

These incidents share a feature: the first line of defense was an operator, not an information technology department. At the same time, roughly one third of water sector workers are nearing retirement age, and EPA has identified cybersecurity among the new skills the replacement workforce will need \cite{epa2026}. What states require operators to learn is therefore a practical question for system resilience, not only an administrative one.

States have begun to respond, but mostly at the level of the utility: risk assessments, cybersecurity programs, incident reporting, and funding. Indiana, for example, has required annual cybersecurity vulnerability assessments of community water systems serving 500 or more people since July 2025 \cite{indiana}, and New York's 2026 rules require covered community water systems to adopt cybersecurity programs from 1 January 2027 \cite{nydoh}. Whether states have also changed what an operator must know to hold a license has not been tracked anywhere. EPA's last national summary of state operator certification programs dates from March 2016 \cite{epa2016}, and it predates every provision described here.

This paper introduces CERT-SCAN, a dataset designed to close that gap. Its contributions are:
\begin{enumerate}
\item A complete, citable baseline of cybersecurity provisions in operator certification for all 50 states, the District of Columbia, and Puerto Rico, coded separately for drinking water and wastewater, with the source behind every record.
\item A five-level coding scheme that separates codified requirements from credit eligibility, general security training, and agency practice, so that the type of legal instrument is visible rather than assumed.
\item A systematic reading of the 2025 national exam content outlines, which shows that cybersecurity has entered certification testing ahead of certification training.
\item An open, scripted pipeline that rebuilds the dataset, figures, and this paper from archived sources, so that the record can be maintained annually.
\end{enumerate}

Section~\ref{sec:background} describes how operator certification works and the 2016 baseline. Section~\ref{sec:methods} sets out the data and methods. Section~\ref{sec:results} reports the findings. Section~\ref{sec:discussion} discusses their implications, Section~\ref{sec:validation} the technical validation, and Section~\ref{sec:limits} the limitations. Section~\ref{sec:conclusion} concludes.

\section{Background}\label{sec:background}

\subsection{How operator certification works}

Operator certification programs share a common structure. A state classifies water and wastewater facilities by size and treatment complexity, assigns operator classes to match, and requires each operator in responsible charge to hold a certificate at or above the facility's class. Certificates are obtained by meeting education and experience requirements and passing a written examination, and they are renewed on a cycle of typically two to five years by completing a set number of approved continuing education hours. Most states administer drinking water certification through a health or environmental agency, sometimes with a board; wastewater certification is frequently run by a different agency or board, under a different statute.

Three levers can bring cybersecurity into this structure. First, a state can require cybersecurity training as a condition of initial certification or renewal. Second, it can permit cybersecurity courses to count toward the continuing education hours operators must earn anyway. Third, the examination can test cybersecurity knowledge. The first two are set by state law, rule, or agency practice. The third is largely set by the exam provider: many states purchase standardized examinations from Water Professionals International (WPI), formerly the examination service of the Association of Boards of Certification, whose need-to-know criteria define what each exam covers \cite{wpimembers,wpintk}.

\subsection{The 2016 national baseline}

EPA's \emph{Summary of State Operator Certification Programs} (EPA 816-R-16-002, March 2016, \epaPages\ pages) profiles the drinking water program of every state and Puerto Rico, two pages each, covering authorization, classification, qualifications, enforcement, renewal, resources, stakeholder involvement, and program review \cite{epa2016}. We searched its full text by script for ten security-related terms. The word ``cyber'' appears \epaCyber\ times. ``Security'' appears \epaSecurity\ time, in Pennsylvania's profile, which notes that every certified operator must complete the Department's system security course. ``SCADA'' appears \epaScada\ time, in a Missouri footnote about operator oversight at facilities with extensive automation. The remaining terms (information technology, operational technology, network, hacker, malware, ransomware) do not appear. The report covers drinking water only; there is no comparable national summary for wastewater.

\section{Data and Methods}\label{sec:methods}

\subsection{Scope and unit of observation}

The unit of observation is a jurisdiction by program: the 50 states, the District of Columbia, and Puerto Rico, each recorded once for drinking water and once for wastewater, for \nRows\ records. Where a single body certifies both groups under one rule, both records carry the same citation, so that counts by program always total \nJuris. All sources were accessed on \accessDate\ unless a record states otherwise.

\subsection{Coding scheme}

Each record is assigned one of five status levels (Table~\ref{tab:scheme}). The level is determined by codified or official instruments: a statute, a regulation, official agency guidance, or a published course approval. The decision rule that most shapes the results is that agency practice without a written rule cannot by itself make a record ``required''; practice that goes beyond the written rule is recorded separately. This rule was fixed by the first author before coding began.

\begin{table}[t]
\caption{CERT-SCAN status levels.}\label{tab:scheme}
\small\centering
\begin{tabularx}{\columnwidth}{@{}>{\raggedright\arraybackslash}p{0.27\columnwidth}>{\raggedright\arraybackslash}X@{}}
\toprule\toprule
Level & Test applied \\
\midrule
Required & A statute or regulation makes cybersecurity-specific training, coursework, or demonstrated knowledge a condition of obtaining or renewing certification. \\
Credited, not required & A rule, official guidance, or published course approval names cybersecurity as eligible for certification credit. \\
General security only & A codified security-training condition exists, but its cybersecurity content is not established. \\
No provision found & No provision was identified in the sources reviewed. A statement about the evidence, not a claim of absence. \\
Unknown & Sources could not be read or conflicted. \\
\bottomrule
\end{tabularx}
\end{table}

For every record, CERT-SCAN also captures the certifying body, the governing statute and rule chapter, the mechanism (new and renewal training, mandatory renewal hours, one-time course, or credit eligibility), hours, renewal cycle, operator classes covered, effective date, instrument type (statute, regulation, emergency rule, guidance, board practice), the exact citation and a link to it, the exam provider, whether the exam outline lists a cybersecurity task, a short paraphrase of the evidence, the sources searched, the access date, a confidence level, and whether an author has verified the record against its source. The codebook in the repository defines every column.

\subsection{Source protocol}

Every record rests on at least one state-specific source. In order of preference this was the renewal or continuing education rule, the certifying body's program page or renewal form, or its training guidance. Sources were located by web search and, where possible, downloaded and logged with their URL, access date, size, and SHA-256 fingerprint. Downloaded text was scanned by script for cybersecurity terms (cyber, SCADA, control system, operational technology, ransomware, malware, network security, information security), and every page with a match was read to judge it. Pages that blocked automated access, or returned challenge pages, were read through a web reader and recorded as such. A page that could not be read was never treated as evidence of absence.

Nine nationwide searches supplemented the state sources, to find statutes, rules, bills, board notices, and course approvals that might not appear in a state's rule chapter. This step found Maryland's 2025 statute, which sits in the state's environment article rather than its certification rules, and identified pending or failed measures in Illinois, Texas, and Utah \cite{bondbuyer,txsb1034,utah}. All queries are listed in the repository.

\subsection{Exam content}

The exam provider for each record was identified from WPI's list of regulatory members, which names the certifying bodies and the exam programs each uses \cite{wpimembers}. WPI's 2025 need-to-know criteria for water treatment, water distribution, wastewater treatment, and wastewater collection, classes 1 to 4 (sixteen outlines), were downloaded and searched for cybersecurity tasks \cite{wpintk}. Membership shows that a jurisdiction uses WPI examinations; it does not show that it has adopted the 2025 version, and the exam column is worded accordingly.

\subsection{Verification and quality assurance}

The records that carry the principal findings (Maryland, New York, and Pennsylvania) were verified by the first author against their primary legal texts, including the chaptered Maryland law, the adopted New York rules and their adoption dates, and the Pennsylvania Bulletin notice that made Pennsylvania's rule effective. Every record states its confidence: high where the primary legal text was read in full, medium where an official agency document was read, and low where the source awaits confirmation. A build script checks every record against the coding rules on each run (Section~\ref{sec:validation}).

\section{Results}\label{sec:results}

\subsection{National overview}

Figure~\ref{fig:dash} summarizes the record. Of the \nJuris\ jurisdictions, \nAnyJurisW\ have some cybersecurity provision in operator certification and \nNoneJuris\ have none in the sources reviewed. By program, drinking water has \dwReq\ required, \dwCred\ credited, \dwGen\ general security, and \dwNone\ no-provision records; wastewater has \wwReq, \wwCred, \wwGen, and \wwNone\ (Fig.~\ref{fig:status}).

\begin{figure*}[t]
\centering
\includegraphics[width=0.94\textwidth]{pbi_fig1_dashboard.png}
\caption{CERT-SCAN overview as of \accessDate. Top: headline counts. Bottom: status of each jurisdiction's drinking water and wastewater certification, on a tile grid of the 50 states, the District of Columbia, and Puerto Rico. A dashed outline marks records awaiting agency confirmation.}
\label{fig:dash}
\end{figure*}

\begin{figure}[t]
\centering
\includegraphics[width=\figw]{pbi_fig2_status_by_program.png}
\caption{Records by status and program. Each program totals \nJuris\ jurisdictions.}
\label{fig:status}
\end{figure}

\subsection{Required provisions}

\nReqJurisCap\ jurisdictions, \reqList, require cybersecurity knowledge for certification by statute or regulation, and both do so for drinking water and wastewater alike (Table~\ref{tab:required}).

\begin{table*}[t]
\caption{Codified cybersecurity requirements in operator certification. DW: drinking water; WW: wastewater. Hours are per renewal cycle.}\label{tab:required}
\small\centering
\begin{tabularx}{\textwidth}{@{}l c >{\raggedright\arraybackslash}p{0.16\textwidth} l c l l >{\raggedright\arraybackslash}X@{}}
\toprule\toprule
Jurisdiction & Program & Mechanism & Hours & Cycle (yr) & Effective & Instrument & Citation \\
\midrule
\tabRequired
\bottomrule
\end{tabularx}
\end{table*}

\emph{Maryland} enacted the first statutory requirement. Chapter 495 of the 2025 Laws of Maryland (Senate Bill 871), approved on 13 May 2025 and effective 1 October 2025, added a subtitle on community water and sewerage system cybersecurity to the Environment Article. Section 9-2702(2) directs the Department of the Environment to include cybersecurity awareness components for all new and renewing operator and superintendent certifications under the state's certification title \cite{mdch495}. The operator board had moved first in practice: from July 2025, renewing water, wastewater, and industrial operators were required to complete EPA's \emph{Cybersecurity 101} training, consisting of a webinar, a video, and a quiz of about thirty minutes, with completion confirmed on the renewal form \cite{mdboard,mdtraining}. Maryland renews certificates every three years. The statute directs the agency rather than binding operators directly, and we did not locate a corresponding amendment to the state's certification regulations; Maryland is nonetheless coded ``required'' because the instrument establishing the requirement is a statute.

\emph{New York} followed with detailed regulations in both programs. For drinking water, Appendix 5-E of Part 5 of the State Sanitary Code was adopted on 11 March 2026. Section 5-E.7, effective on adoption, requires every operator certified under Subpart 5-4 to complete at least one hour of Department-approved cybersecurity training every three years, for new and renewal certifications alike, with no exclusions \cite{nyapp5e,nydoh}. For wastewater, amendments to 6 NYCRR Part 650, effective 26 March 2026, require cybersecurity training for renewal: two contact hours for grades 1/1A and 2/2A and four for grades 3/3A and 4/4A per five-year renewal period, counted within the existing continuing education total, for certificates expiring on or after 1 January 2027, using Department-approved courses only \cite{nydec}.

The two states' designs differ in ways a state drafting its own rule would need to decide: half an hour versus one to four hours; a three-year versus a five-year cycle; a federal course delivered as is versus state-approved curricula; and, in New York's wastewater rule, hours carved out of the existing continuing education total rather than added to it.

\subsection{Credited provisions}

\nCredJurisCap\ jurisdictions credit cybersecurity training toward certification without requiring it (Table~\ref{tab:credited}). The mechanisms differ. In \emph{Texas}, every water operator has since 1 April 2024 had to complete a two-hour resiliency overview for a new or upgraded license and two hours of approved resiliency training in each renewal cycle; cybersecurity is a permitted element of one of the required topics, but not required itself, and the requirement does not extend to wastewater operators \cite{tceq}. \emph{Montana}'s continuing education guidelines give full credit to cybersecurity courses for water and wastewater operators and list cybersecurity among the approved safety topics \cite{montana}. \emph{Maine}'s wastewater renewal guidance lists cybersecurity as an acceptable operations and maintenance topic \cite{maine}. \emph{Michigan}'s environmental agency approved its cybersecurity webinar series for both drinking water and wastewater continuing education credit \cite{michigan}. \emph{Massachusetts}' drinking water board approved six cybersecurity courses for training contact hours in March 2026, and the state offers a free self-paced course worth one contact hour \cite{massboard,massnews}.

\begin{table*}[t]
\caption{Provisions that credit, but do not require, cybersecurity training.}\label{tab:credited}
\small\centering
\begin{tabularx}{\textwidth}{@{}l c >{\raggedright\arraybackslash}p{0.2\textwidth} l l >{\raggedright\arraybackslash}X@{}}
\toprule\toprule
Jurisdiction & Program & Mechanism & Instrument & Date & Source \\
\midrule
\tabCredited
\bottomrule
\end{tabularx}
\end{table*}

The credited category is almost certainly undercounted. Certifying bodies approve individual courses routinely, and most approvals are never published; only jurisdictions where an official document names cybersecurity, or a published approval exists, are coded as credited.

\subsection{General security training}

\emph{Pennsylvania} is the only jurisdiction with a codified security-training condition predating 2025. Under 25 Pa.\ Code \S\,302.804, every certified operator must complete a Department-approved system security training course; the rule was published and took effect on 18 September 2010 (40 Pa.B.\ 5294) \cite{pa302}. The course, \emph{Securing Drinking Water and Wastewater Treatment Facilities}, carries 3.5 contact hours, is taken once, and failing to complete it results in loss of the license; the Department may also require further security courses by written notice \cite{pa302,pasec01}. Its published description does not state whether it covers cybersecurity, and confirmation has been requested from the Department. If it does, Pennsylvania's requirement predates every other provision in this record by fifteen years.

\subsection{No provision found}

In the remaining records (\dwNone\ drinking water and \wwNone\ wastewater), no cybersecurity provision was found in the state's renewal rule, program page, or guidance. Several of these sources are detailed. Ohio's renewal rule, revised in July 2025, requires operations and maintenance hours and a rule-review course but names no security topic \cite{ohio}. Florida's renewal rule lists ``security, and emergency response'' among eligible continuing education topics without naming cybersecurity \cite{florida}. Three automated matches proved not to be provisions: North Carolina's wastewater commission is named for ``water pollution control system'' operators; Vermont's wastewater rule awards facility classification points for a computer-based control system; and Wyoming's program page lists a federal cybersecurity advisor as a contact. The District of Columbia is coded with low confidence because no independent District certification program for water or wastewater operators was identified; a 2016 DC Water document describes certification as a requirement outside Washington, DC \cite{dcwater}, and confirmation from the District has been requested.

\subsection{Testing ahead of training}

The exam outlines tell a different story from the rules (Table~\ref{tab:exam}, Fig.~\ref{fig:exam}). In WPI's 2025 need-to-know criteria, all four water treatment classes list ``perform cyber security according to industry standards'' among the tasks in the \emph{Security, Safety, and Administrative Procedures} domain. That domain accounts for 12 to 14 of the 100 scored questions, depending on class, and covers roughly twenty tasks, so cybersecurity is in scope but lightly weighted. None of the distribution, wastewater treatment, or collection outlines mentions cybersecurity \cite{wpintk}.

\begin{table}[t]
\caption{Cybersecurity task in the 2025 WPI need-to-know criteria, by exam and class.}\label{tab:exam}
\small\centering
\begin{tabular}{@{}lcccc@{}}
\toprule\toprule
Exam & Class 1 & Class 2 & Class 3 & Class 4 \\
\midrule
\tabExam
\bottomrule
\end{tabular}
\end{table}

WPI examinations are used for \nWPIdw\ drinking water programs whose water treatment exams therefore include a cybersecurity task, and for \nWPIww\ wastewater programs whose exams do not (Fig.~\ref{fig:coverage}); \nOtherExamDW\ drinking water and \nOtherExamWW\ wastewater programs use other or state-developed examinations. Of the \nWPIdw\ drinking water programs using these water treatment exams, \nWPIdwNoReq\ are in jurisdictions with no training requirement. If those programs have adopted the 2025 outlines, most water treatment operators who may now be tested on cybersecurity work in states that do not require them to be trained in it.

\begin{figure}[t]
\centering
\includegraphics[width=\figw]{pbi_fig4_exam_matrix.png}
\caption{Cybersecurity task in the 2025 WPI exam outlines, by exam area and class.}
\label{fig:exam}
\end{figure}

\begin{figure}[t]
\centering
\includegraphics[width=\figw]{pbi_fig5_exam_coverage.png}
\caption{Programs by exam source and whether the exam outline lists a cybersecurity task.}
\label{fig:coverage}
\end{figure}

\subsection{Timing and instruments}

Figure~\ref{fig:timeline} places each provision on a timeline. Apart from Pennsylvania's 2010 rule, every provision dates from 2022 or later, and both codified requirements took effect within six months of each other, between October 2025 and March 2026. The instruments are mixed (Fig.~\ref{fig:instrument}): of the \nPositive\ records with a provision, \nInstStatute\ rest on a statute, \nInstReg\ on regulations, \nInstGuidance\ on agency guidance, and \nInstBoard\ on board course approvals.

\begin{figure*}[t]
\centering
\includegraphics[width=0.94\textwidth]{pbi_fig3_timeline.png}
\caption{Effective date of each cybersecurity provision in operator certification, by status. DW: drinking water; WW: wastewater. Dates for credited provisions are the dates of the guidance or approval cited.}
\label{fig:timeline}
\end{figure*}

\begin{figure}[t]
\centering
\includegraphics[width=\figw]{pbi_fig6_instrument.png}
\caption{Instrument type behind each record with a provision.}
\label{fig:instrument}
\end{figure}

\section{Discussion}\label{sec:discussion}

\subsection{A testing and training gap}

The clearest finding is a mismatch between what operators may be examined on and what they are required to learn. The water treatment exams used by most states now list a cybersecurity task, yet only two states require training. Two readings are possible, and the data do not decide between them. One is that the exam provider is ahead of regulators and states will follow, as Utah's environmental agency is now considering \cite{utah}. The other is that a single task in a lightly weighted domain will have little effect on what operators learn unless renewal training reinforces it. Either way, the gap is measurable, and the annual update will show whether it closes.

The gap is also uneven across the workforce. Distribution, wastewater treatment, and collection operators are not tested on cybersecurity, and wastewater has fewer credited provisions than drinking water. Yet the summer 2026 attacks targeted wastewater as well as drinking water utilities \cite{axios2026}.

\subsection{What a requirement looks like}

The two codified requirements show the design choices a state faces: whether to rely on a federal course or approve its own curricula, how many hours to require, whether to apply the requirement at initial certification, at renewal, or both, and whether to add hours or carve them out of existing totals. Maryland's model is light and fast: one short federal course, implemented by the board before the statute took effect. New York's is heavier and more tailored: approved curricula, hours scaled to wastewater grade, and coordination with utility-level cybersecurity programs. CERT-SCAN records these parameters so that states drafting provisions can compare options directly.

\subsection{Why the instrument matters}

Of the \nPositive\ records with a provision, \nNonCodified\ rest on guidance or course approvals rather than on statute or regulation. These are faster to adopt and easier to change, but also easier to miss: Maryland's requirement would have been invisible to a review limited to certification rules, because it sits in a separate statute. A record that names the instrument behind each provision lets users distinguish a durable legal requirement from a policy that a new board could reverse.

\subsection{Uses of the record}

We see four immediate uses. State programs can locate peers and model language. EPA and associations of state program managers can track adoption and target technical assistance. Trainers and exam developers can align course content with the provisions that award credit. Researchers can treat the record as a policy variable, linking it in future to incident, compliance, or workforce outcomes, which CERT-SCAN does not itself measure.

\section{Technical Validation}\label{sec:validation}

The build script checks, on every run, that each jurisdiction and program appears exactly once; that every status and instrument value is valid; that no record is coded ``required'' without a statute, regulation, or emergency rule; that every record with a provision has a source link; and that every record without one documents the sources searched. The current build reports \nQA\ problems.

Of the \nRows\ records, \nStateSpecific\ rest on a state-specific source: \nHigh\ on primary legal text read in full and \nMedium\ on an official agency document. \nLow\ records (the District of Columbia) await confirmation. The first author has verified \nVerified\ records against their sources: every record with a provision, and a random sample of \nSample\ records coded ``no provision found,'' drawn with a fixed seed so that the sample can be reproduced; all sampled records matched their sources (Fig.~\ref{fig:evidence}). The 2016 baseline scan, the state-source scan, and the exam outline scan are scripted and reproduce from the archived files, each logged with its SHA-256 fingerprint.

\begin{figure}[t]
\centering
\includegraphics[width=\figw]{pbi_fig7_evidence.png}
\caption{Evidence behind each record: primary legal text (high confidence), official agency document (medium), or awaiting confirmation (low).}
\label{fig:evidence}
\end{figure}

\section{Limitations}\label{sec:limits}

CERT-SCAN measures whether rules require cybersecurity knowledge, not whether operators have it, how good the training is, or whether any system is more secure. A finding of ``no provision found'' reflects the sources reviewed, typically one renewal rule or program page per record together with a nationwide search; a provision in an unread document, or adopted after the access date, would be missed. Because course approvals are often unpublished, the credited category is undercounted. Exam content is inferred from published outlines, not from exam items, and WPI membership does not establish adoption of the 2025 outlines. Pennsylvania's course content and the District of Columbia's program status await confirmation from the responsible agencies. The 2016 EPA baseline covers drinking water only. Effective dates for credited provisions are the dates of the guidance or approval cited, which may be later than when credit was first given. Finally, the record is a point-in-time snapshot of a fast-moving area; Illinois, Utah, and others have measures pending.

\section{Conclusion}\label{sec:conclusion}

Ten years after EPA's last national summary of operator certification programs, which did not mention cybersecurity, \nReqJurisW\ jurisdictions now require it, \nCredJurisW\ credit it, and the national water treatment exams test it. Most jurisdictions have no provision at all, and wastewater lags drinking water on every measure. CERT-SCAN makes this landscape visible, cites every record, and will be updated annually so that change can be tracked rather than inferred. We invite certifying bodies, exam providers, and researchers to correct, extend, and use it.

\section*{Data Availability}
The dataset (\texttt{cert\_scan\_v1.csv}, \nRows\ records), codebook, search log, provenance log, quality-assurance report, and figures are available at \url{\repoURL}. Version 1.0 is archived on Zenodo, \href{https://doi.org/10.5281/zenodo.23608016}{doi:10.5281/zenodo.23608016}. Data and documentation are licensed CC BY 4.0.

\section*{Code Availability}
The scripts that download sources, scan text, build the dataset, draw the figures, and compile this paper are in the same repository under the MIT license.

\section*{Use of AI Tools}
AI assistance (Claude, Anthropic) was used to search for and retrieve sources, write data-processing and figure code, and draft text. The authors made all coding decisions, verified the principal records against primary sources, and take responsibility for the content.

\section*{Author Contributions and Competing Interests}
\authorName: conceptualization, methodology, data curation, investigation, validation, and writing. \coName: \coRoles. The authors declare no competing interests and received no funding for this work.

\section*{Acknowledgment}
The authors thank the state certification programs whose public documents make this record possible.

\begin{thebibliography}{99}\small\setlength{\itemsep}{1pt}
\bibitem{sdwa} Safe Drinking Water Act \S\,1419, 42 U.S.C.\ \S\,300g-8.
\bibitem{epa1999} U.S.\ EPA, Final guidelines for the certification and recertification of the operators of community and nontransient noncommunity public water systems, 64 Fed.\ Reg.\ 5916 (5 Feb.\ 1999).
\bibitem{epaopcert} U.S.\ EPA, Operator certification program fact sheet, \url{https://www.epa.gov/sites/production/files/2015-04/documents/opcert.pdf}.
\bibitem{securityweek} SecurityWeek, ICS at multiple US water facilities targeted by hackers affiliated with Iranian government (4 Dec.\ 2023).
\bibitem{epa2024} U.S.\ EPA, Enforcement alert: drinking water systems to address cybersecurity vulnerabilities (May 2024), \url{https://www.epa.gov/node/289709}.
\bibitem{axios2026} Axios, The number of states targeted in cyberattacks on water systems has jumped to 12 (4 Aug.\ 2026).
\bibitem{epa2026} U.S.\ EPA, EPA convenes roundtable focused on strengthening the water sector workforce (5 Aug.\ 2026).
\bibitem{indiana} Indiana Department of Environmental Management, Drinking water operator certification program report to the Governor (29 Sep.\ 2026).
\bibitem{nydoh} New York State Department of Health, Drinking water system cybersecurity requirements table.
\bibitem{epa2016} U.S.\ EPA, \emph{Summary of State Operator Certification Programs}, EPA 816-R-16-002 (March 2016).
\bibitem{wpimembers} Water Professionals International, WPI members, \url{https://gowpi.org/about-us/wpi-members/} (accessed 23 Oct.\ 2026).
\bibitem{wpintk} Water Professionals International, 2025 need-to-know criteria: water treatment, water distribution, wastewater treatment, and wastewater collection, classes 1--4, \url{https://gowpi.org/services/2025-need-to-know-criteria/}.
\bibitem{bondbuyer} The Bond Buyer, Illinois bill strengthens water utility cybersecurity rules (28 Sep.\ 2026).
\bibitem{txsb1034} Texas Legislature, S.B.\ 1034 (2025), introduced, not enacted.
\bibitem{utah} Utah Legislature, 2026 interim audit follow-up, Department of Environmental Quality response, \url{https://le.utah.gov/interim/2026/pdf/00000993.pdf}.
\bibitem{mdch495} Maryland, 2025 Md.\ Laws ch.\ 495 (S.B.\ 871), Md.\ Code Ann., Envir.\ \S\,9-2702.
\bibitem{mdboard} Maryland Department of the Environment, Board of Waterworks and Waste Systems Operators, meeting minutes (20 Mar.\ 2025) and board page.
\bibitem{mdtraining} Maryland Department of the Environment, Cybersecurity training for water and wastewater system operators.
\bibitem{nyapp5e} New York, 10 NYCRR Part 5, Appendix 5-E, \S\,5-E.7 (adopted 11 Mar.\ 2026).
\bibitem{nydec} New York State Department of Environmental Conservation, 6 NYCRR \S\S\,650.8 and 650.12; Wastewater cybersecurity resources.
\bibitem{pa302} Pennsylvania, 25 Pa.\ Code \S\,302.804; 40 Pa.B.\ 5294 (18 Sep.\ 2010).
\bibitem{pasec01} Pennsylvania DEP, SEC01 Securing Drinking Water and Wastewater Treatment Facilities, course 2954.
\bibitem{tceq} Texas Commission on Environmental Quality, RG-637, Resiliency requirements for water operators and training providers.
\bibitem{montana} Montana DEQ, Continuing education credit guidelines (8 Jun.\ 2023), \S\S\,3.1, 5.2.
\bibitem{maine} Maine DEP, Training contact hour requirements for certification renewal, v1.3 (2022), Table 1.
\bibitem{michigan} Michigan EGLE, Cybersecurity for the water sector webinar series.
\bibitem{massboard} Massachusetts Board of Certification of Operators of Drinking Water Supply Facilities, meeting agenda (13 Mar.\ 2026).
\bibitem{massnews} MassDEP Drinking Water Program, \emph{In the Main} (11 Aug.\ 2026).
\bibitem{ohio} Ohio Adm.\ Code 3745-7-15 (effective 3 Jul.\ 2025).
\bibitem{florida} Fla.\ Admin.\ Code r.\ 62-602.710.
\bibitem{dcwater} DC Water, Human Resources and Labor Relations Committee materials (13 Jan.\ 2016).
\end{thebibliography}

\clearpage
\onecolumn
\appendix
\renewcommand{\thesection}{A}
\section*{Appendix: CERT-SCAN record by jurisdiction}
\small
Status: R = required; C = credited, not required; G = general security only; -- = no provision found. Exam: WPI+ = WPI exam whose 2025 outline lists a cybersecurity task; WPI = WPI exam without one; Other = other or state exam. Evidence: H = primary legal text; M = official agency document; L = awaiting confirmation. Full citations, links, and notes for each record are in the dataset.
\vspace{4pt}
\renewcommand{\arraystretch}{1.08}
\begin{longtable}{@{}>{\raggedright\arraybackslash}p{0.14\textwidth}ccccc>{\raggedright\arraybackslash}p{0.41\textwidth}@{}}
\toprule\toprule
Jurisdiction & DW & WW & Exam DW & Exam WW & Evidence & Program authority \\
\midrule\endfirsthead
\toprule
Jurisdiction & DW & WW & Exam DW & Exam WW & Evidence & Program authority \\
\midrule\endhead
\bottomrule\endfoot
\tabAppendix
\end{longtable}
\end{document}
